\sectionCenteredToc{ПЕРЕЧЕНЬ СОКРАЩЕНИЙ}

% Оборачиваем в sloppy, чтобы разрешить выравнивание по ширине для длинных слов
{\sloppy
% Настраиваем 3 колонки:
% 1. p{0.136\textwidth} - ширина под аббревиатуру (как ты замерял)
% 2. @{} c @{\hspace{0.4em}} - узкая колонка под тире, прижатая к тексту
% 3. p{0.80\textwidth} - колонка для текста (ровный кирпич, выровненный по ширине)
\begin{longtable}{@{} p{0.18\textwidth} @{} c @{\hspace{0.4em}} p{0.75\textwidth} @{}}
ГОСТ & -- & государственный стандарт \\
ЕСПД & -- & единая система программной документации \\
ИСО/МЭК & -- & Международная организация по стандартизации / Международная электротехническая комиссия (ISO/IEC) \\
МПК & -- & международная патентная классификация \\
ПЭВМ & -- & персональная электронно-вычислительная машина \\
СанПиН & -- & санитарные правила и нормы \\
СНБ & -- & строительные нормы Республики Беларусь \\
СТБ & -- & стандарт Республики Беларусь \\
УЗО & -- & устройство защитного отключения \\
ФИПС & -- & Федеральный институт промышленной собственности \\
AI & -- & (Artificial Intelligence) искусственный интеллект \\
API & -- & (Application Programming Interface) описание способов взаимодействия приложений \\
ARM & -- & (Advanced RISC Machine) семейство микропроцессорных архитектур с сокращенным набором команд \\
AVX2 & -- & (Advanced Vector Extensions 2) векторное расширение системы команд x86 для микропроцессоров \\
BGR & -- & (Blue, Green, Red) цветовая модель, аналогичная RGB, но с обратным порядком следования цветовых каналов \\
BYN & -- & (Belarusian Ruble) белорусский рубль, национальная валюта Республики Беларусь \\
CNN & -- & (Convolutional Neural Network) сверточная нейронная сеть \\
CPU & -- & (Central Processing Unit) центральный процессор \\
DCT & -- & (Discrete Cosine Transform) дискретное косинусное преобразование \\
DIV2K & -- & (DIVerse 2K resolution images) набор данных изображений в разрешении 2К \\
DML & -- & (Data manipulation language) язык манипуляции данных \\
DNN & -- & (Deep Neural Network) глубокая нейронная сеть \\
DWT & -- & (Discrete Wavelet Transform) дискретное вейвлет-преобразование \\
Edge Computing & -- & периферийные вычисления, парадигма распределенных вычислений \\
Edge-to-Cloud & -- & модель распределенных вычислений с переносом ресурсоемких задач в облако \\
EDSR & -- & (Enhanced Deep Residual Networks) улучшенная глубокая остаточная нейронная сеть \\
ESPCN & -- & (Efficient Sub-Pixel Convolutional Neural Network) эффективная субпиксельная сверточная нейронная сеть \\
FSRCNN & -- & (Fast Super-Resolution Convolutional Neural Network) быстрая сверточная нейронная сеть \\
GNU & -- & (GNU's Not Unix) свободная UNIX-подобная операционная система \\
GPU & -- & (Graphics Processing Unit) графический процессор \\
GUI & -- & (Graphical User Interface) графический пользовательский интерфейс \\
HR & -- & (High Resolution) высокое разрешение \\
HTTP & -- & (HyperText Transfer Protocol) протокол прикладного уровня \\
IoT & -- & (Internet of Things) интернет вещей \\
JPEG & -- & (Joint Photographic Experts Group) популярный графический формат сжатия с потерями \\
LED & -- & (Light-Emitting Diode) светоизлучающий диод \\
LL & -- & (Low-Low) низкочастотная субполоса аппроксимации вейвлет-преобразования \\
LR & -- & (Low Resolution) низкое разрешение \\
MAE & -- & (Mean Absolute Error) средняя абсолютная ошибка \\
MSE & -- & (Mean Squared Error) среднеквадратичная ошибка \\
NumPy & -- & (Numerical Python) фундаментальная библиотека для языка Python \\
OOM & -- & (Out-Of-Memory) состояние исчерпания оперативной памяти \\
OpenCV & -- & (Open Source Computer Vision Library) открытая библиотека алгоритмов компьютерного зрения \\
POST & -- & метод запроса протокола HTTP для отправки объемных данных \\
PS & -- & (PixelShuffle) операция субпиксельной свертки для увеличения разрешения \\
PSNR & -- & (Peak Signal-to-Noise Ratio) пиковое отношение сигнала к шуму \\
Python & -- & высокоуровневый язык программирования общего назначения \\
RAM & -- & (Random Access Memory) оперативная память \\
ReLU & -- & (Rectified Linear Unit) нелинейная функция активации искусственного нейрона \\
REST & -- & (Representational State Transfer) архитектурный стиль взаимодействия компонентов распределенного приложения в сети \\
RGB & -- & (Red, Green, Blue) аддитивная цветовая модель \\
SISR & -- & (Single Image Super-Resolution) восстановление сверхразрешения одного изображения \\
SOAP & -- & (Simple Object Access Protocol) протокол обмена структурированными сообщениями в распределенной вычислительной среде \\
SOQL & -- & (Salesforce object query language) язык запросов к базе данных \\
SSD & -- & (Solid-State Drive) твердотельный накопитель \\
SSIM & -- & (Structural Similarity Index Measure) индекс структурного сходства \\
VMM & -- & (Virtual machine monitor) монитор виртуальных машин \\
VRAM & -- & (Video Random Access Memory) видеопамять графического ускорителя \\
WAN & -- & (Wide area network) глобальная сеть \\
XML & -- & (Extensible Markup Language) расширяемый язык разметки \\
YCbCr & -- & цветовая модель (базовая яркость и две цветоразностные компоненты) \\
\end{longtable}
}